\subsection{A dyadic majorant with an exact sum}

The quantum experiment gives the probabilities in
\eqref{eq:residue-probability}, but their dependence on the bin is inconvenient.
We will supplement the experiment with classical random choices so that, at
the true order \(d=r\), every reduced numerator contributes exactly
\(\lambda(d)\).  The next lemma supplies the arithmetic needed to do this
without knowing \(r\). The construction has three steps: sum each
progression exactly, approximate its mass by a piecewise polynomial, and
round that polynomial upward so that its total is computable without
enumerating residues. Mosca and Zalka use power sums to approximate a
success probability to arbitrary precision~\cite[Section~4.1]{MZ}.
Here the polynomial pieces have bounded degree, so their dyadic envelope
can be summed exactly.

\begin{lemma}\label{lem:majorant}
For every \(1\le d<B\) and \(0\le j<d\) with \(\gcd(j,d)=1\), there are
nonnegative dyadic numbers \(U(d,j,u)\), \(0\le u<d\), such that
\begin{align}
 0&\le U(d,j,jt\bmod d)-P(d,j,t)\le \frac{128}{Q}
       &&(0\le t<d), \label{majorant-error}\\
 \frac{1}{128d}&\le
 S(d,j):=\frac12\sum_{u=0}^{d-1}U(d,j,u)\le2.
 \label{majorant-sum}
\end{align}
Given their integer arguments, \(P(d,j,t)\), \(U(d,j,u)\), and \(S(d,j)\)
are exactly computable in time polynomial in \(n\), and are dyadic numbers
of polynomial bit length.  In particular, computing \(S(d,j)\) does not
require enumerating its \(d\) summands.
\end{lemma}

\begin{proof}
Fix \(d,j\), and abbreviate \(k=k(d,j)\) and \(\eta=k-jQ/d\).
For \(x\equiv t\pmod d\), put \(u=jt\bmod d\), represented in
\(\{0,\ldots,d-1\}\).  Periodicity gives the exact identity
\[
 g(kx/Q)=g(u/d+\eta x/Q).
\]
As \(x\) traverses the progression, this last argument moves a distance at
most \(1/2\).  Thus it meets only a bounded number of the quarter-integer
breakpoints of \(g\).  Split the progression at those breakpoints by rational
comparisons and integer division.  On each resulting part, the real and
imaginary coordinates are affine in the progression index, so their sums
are obtained from the number of terms and the sum of consecutive integers.
This computes \(P(d,j,t)\) with a bounded number of arithmetic operations
on \(O(n)\)-bit integers and rationals.  Its dyadic nature is also immediate
from the original sum: each \(g(kx/Q)\) has real and imaginary denominators
dividing \(Q\), so the squared modulus after division by \(Q\) has denominator
dividing \(Q^4\).

Although each \(P(d,j,t)\) is efficiently computable, direct summation over
\(t\) requires \(d\) evaluations, which need not be polynomial in \(n\).
To construct a majorant with an efficiently computable total, first define,
for real \(0\le u<d\),
\begin{equation}\label{majorant-integral}
 I(u)=\frac1d\int_0^1g(u/d+\eta v)\,dv,
 \qquad J(u)=|I(u)|^2.
\end{equation}
We first record two uniform estimates:
\begin{equation}\label{majorant-integral-bounds}
 |P(d,j,t)-J(jt\bmod d)|\le\frac{32}{Q},
 \qquad
 \frac1{8d}\le |I(u)|\le\frac1d.
\end{equation}
For the first estimate, let \(h=d/Q<1\) and
\(f(v)=g(u/d+\eta v)\).  The function \(g\) has magnitude at most one and
Lipschitz constant \(4\sqrt2\), so \(f\) has Lipschitz constant less than
four.  The progression points \(x/Q\) are the left endpoints of intervals
of length \(h\).  Together these intervals start within \(h\) of zero and
end between \(1\) and \(1+h\).  Their total length is at most \(1+h\).
The error of the left rectangle rule is at most \(2h(1+h)\), and the two
endpoint discrepancies contribute at most \(2h\).  Consequently, writing
\[
 A_t=\frac1Q\sum_{\substack{0\le x<Q\\x\equiv t\;(\mathrm{mod}\ d)}}g(kx/Q),
\]
we have \(|A_t-I(u)|\le 6/Q\le8/Q\).  Since
\(|A_t|\le 1/d+1/Q\) and \(|I(u)|\le1/d\), this implies the first
inequality in \eqref{majorant-integral-bounds}.

For the lower bound, replace \(g\) in \eqref{majorant-integral} by
\(e^{-2\pi i z}\).  The resulting integral has magnitude
\[
 \frac1d\left|\frac{\sin(\pi\eta)}{\pi\eta}\right|
 \ge \frac2{\pi d},
\]
where the quotient is interpreted as one at \(\eta=0\).
The inequality follows from concavity of sine on \([0,\pi/2]\).
Linear interpolation on intervals of length \(1/4\) changes the
exponential by at most
\[
 \frac18\left(\frac14\right)^2
 \sup_z\left|\frac{d^2}{dz^2}e^{-2\pi i z}\right|
 =\frac{\pi^2}{32}.
\]
This interpolation bound follows by Taylor expansion at the interpolated
point and linear averaging of the two endpoint remainders; it applies
equally to complex-valued functions.  Therefore
\[
 |I(u)|\ge \frac1d\left(\frac2\pi-\frac{\pi^2}{32}\right)>\frac1{8d}.
\]
For completeness, \(\pi<2\sqrt3\), obtained from
\(\tan(\pi/6)>\pi/6\), bounds the parenthesis below by
\(1/\sqrt3-3/8>1/8\).  The upper bound follows from \(|g|\le1\).

It remains to construct a dyadic majorant whose sum is accessible.
Let \(G\) be the continuous primitive of \(g\) satisfying \(G(0)=0\).
Its real and imaginary parts are quadratic polynomials with rational
coefficients on each quarter interval.  When \(\eta\ne0\),
\begin{equation}\label{majorant-primitive}
 I(u)=\frac{G(u/d+\eta)-G(u/d)}{d\eta}.
\end{equation}
The two arguments range in \([-1/2,3/2]\), so only a bounded number of
pieces occur.  Their boundaries in \(u\) arise when either argument is a
quarter-integer.  On each piece, \(I\) is a complex polynomial in \(u\) of
degree at most two; hence \(J\) is a rational polynomial of degree at most
four.  If \(\eta=0\), the identity \(I(u)=g(u/d)/d\) gives the same conclusion.
All boundaries and coefficients have \(O(n)\) bits and can be computed with
a bounded number of rational operations.  In particular,
\(\eta=(kd-jQ)/d\), so every nonzero \(\eta\) has
\(|\eta|\ge1/d\); the division in \eqref{majorant-primitive} has no hidden
precision cost.

On each piece write \(J(u)=\sum_{q=0}^4a_qu^q\), allowing zero coefficients,
and set \(D=QB^8\).  Replace \(a_q\) by
\(\widetilde a_q=\lceil Da_q\rceil/D\).  Denote the resulting piecewise
polynomial by \(\widetilde J\), assigning boundary points consistently to
one adjacent piece, and define
\begin{equation}\label{majorant-construction}
 U(d,j,u)=\widetilde J(u)+\frac{64}{Q}
 \qquad (u=0,\ldots,d-1).
\end{equation}
Because \(u\ge0\), rounding each coefficient upward increases the value
even when that coefficient is negative.  Moreover,
\[
 0\le \widetilde J(u)-J(u)
 \le \frac{1+u+u^2+u^3+u^4}{QB^8}\le\frac1Q.
\]
Together with \eqref{majorant-integral-bounds}, this yields
\(32/Q\le U(d,j,jt\bmod d)-P(d,j,t)\le97/Q\), proving
\eqref{majorant-error}.  It also gives
\[
 \frac1{128d}
 \le\frac12\sum_{u=0}^{d-1}U(d,j,u)
 \le\frac1{2d}+\frac{65d}{2Q}\le2,
\]
where the last inequality uses \(d<B\) and \(Q=B^{16}\).

Finally, intersect the bounded collection of rational piece intervals with
the integers \(0,\ldots,d-1\), using floors and ceilings and assigning each
boundary integer only once.  Sum each polynomial by the integer power sums
through degree four.  These sums are computed, for example, recursively
from
\[
 M^{q+1}=\sum_{\ell=0}^q\binom{q+1}{\ell}
             \sum_{u=0}^{M-1}u^\ell
 \qquad (0\le q\le4).
\]
Here the zeroth power sum is the count \(M\), with \(u^0=1\) also at zero.
These sums have \(O(n)\) bits for \(M\le B\), and there are only a bounded number
of pieces.  This computes \(S(d,j)\) in polynomial time without enumerating
the residues.  Since \(D\) is a power of two, both \(U\) and \(S\) are
dyadic with \(O(n)\)-bit representations.  All operations used to compute
them are rational arithmetic, integer division, and rounding.
\end{proof}

\subsection{Three modes with a prescribed mass}

Lemma~\ref{lem:majorant} lets us add the discrepancy between \(P\) and \(U\)
by uniform random choices.  A further, smaller correction will make the
mass per reduced numerator exactly \(\lambda(d)\).

\begin{proof}[Completion of the proof of Proposition~\ref{prop:order-trial}]
For each reduced label, the natural scaling coefficient
\(\lambda(d)/S(d,j)\) need not be dyadic. We round it downward and reserve
the remaining mass for the third mode. Define
\begin{equation}\label{mode-coefficient}
 c(d,j)=2^{-10n}
   \left\lfloor 2^{10n}\frac{\lambda(d)}{S(d,j)}\right\rfloor,
 \qquad
 \rho(d,j)=\lambda(d)-c(d,j)S(d,j).
\end{equation}
By \eqref{majorant-sum} and \(\lambda(d)\le1/(1024d)\),
\begin{equation}\label{mode-coefficient-bounds}
 0\le c(d,j)\le\frac18,\qquad
 0\le\rho(d,j)\le2\cdot2^{-10n}.
\end{equation}
Both numbers are exactly computable dyadics of polynomial bit length.

Choose one of the following three modes with probabilities \(1/2,1/4,1/4\),
respectively.  In every mode, reject unless
\[
 1\le d<B,\qquad 0\le j<d,\qquad \gcd(j,d)=1,
 \qquad a^d\equiv1\pmod m.
\]
Check the label conditions before evaluating any formula involving it.
An invalid label or a failed retention test produces no candidate.
\begin{enumerate}
 \item Run the quantum experiment.  Require its phase-register success,
 recover the reduced label \((d,j)\) from the observed \(k\), and retain
 the candidate \(d\) with probability \(c(d,j)\).
 \item Choose independent uniform \(n\)-bit labels \(d,j,t\).
 Require also \(t<d\), and retain \(d\) with probability
 \begin{equation}\label{mode-discrepancy}
  2B^3\bigl(U(d,j,jt\bmod d)-P(d,j,t)\bigr)c(d,j).
 \end{equation}
 \item Choose independent uniform \(n\)-bit labels \(d,j\), and retain
 \(d\) with probability
 \begin{equation}\label{mode-remainder}
  4B^2\rho(d,j).
 \end{equation}
\end{enumerate}
All retention probabilities are dyadic and belong to \([0,1]\).
Nonnegativity follows from \eqref{majorant-error} and
\eqref{mode-coefficient-bounds}; the last two probabilities are at most
\[
 2B^3\frac{128}{Q}\le256B^{-13}\le1,
 \qquad
 4B^2\cdot2\cdot2^{-10n}=8B^{-8}\le1.
\]
Each retention test can therefore be implemented exactly by comparing
fresh fair bits with an integer numerator.  The number of bits and all
arithmetic and quantum costs are polynomial in \(n\).

To compute the mass of the correct order, fix \(d=r\) and a reduced
numerator \(j\).  By \eqref{eq:residue-probability}, Mode 1 contributes
\[
 \frac{c(d,j)}2\sum_{t=0}^{d-1}P(d,j,t).
\]
Each triple in Mode 2 has probability \(B^{-3}\) conditional on choosing
that mode.  Thus Mode 2 contributes
\[
 \frac{c(d,j)}2\sum_{t=0}^{d-1}
       \bigl(U(d,j,jt\bmod d)-P(d,j,t)\bigr).
\]
Multiplication by \(j\) permutes the residues modulo \(d\), so these two
contributions sum to \(c(d,j)S(d,j)\).  Mode 3 contributes
\(\rho(d,j)\), making the total exactly \(\lambda(d)\).
The bin labels are distinct by Lemma~\ref{lem:bins}, and there are
\(\varphi(r)\) reduced numerators.  Hence
\[
 \Pr(\text{the returned candidate is }r)=\varphi(r)\lambda(r).
\]
This includes \(r=1\), whose sole reduced label is \(j=0\).

Every returned candidate satisfies \(a^d\equiv1\pmod m\), so it is a
multiple of the true order \(r\).  At false labels no identity between
the artificial quantity \(P(d,j,t)\) and a quantum probability was used.
This proves all claims of Proposition~\ref{prop:order-trial}.
\end{proof}
